\section{Выводы}

Робот написан и работает: обходит два источника, сохраняет документы в базу,
переживает остановку и умеет обновлять то, что изменилось. Переобкачка вынесена
в отдельную программу и связана с роботом через очередь.

Самое полезное, что дала работа, --- это несколько ошибок, которые на бумаге не
видны.

\textbf{Фильтр \enquote{что обходить} и фильтр \enquote{что сохранять} --- это
разные фильтры.} В первой версии робот складывал в базу оглавления и каталоги:
текста в них достаточно, чтобы пройти проверку по длине, и формально они
проходили все условия. Понадобился отдельный набор правил, решающий, что
считать документом, а что --- только источником ссылок.

\textbf{Без ограничения на то, откуда брать ссылки, корпус перестаёт быть
тематическим.} Из статьи о фильме ссылки ведут на статьи о городах, годах и
профессиях, и через несколько шагов обхода робот собирает уже не кино. Пришлось
разрешить сбор ссылок только со страниц категорий и каталогов.

\textbf{Стандартная библиотека умеет подвести молча.} Робот отказывался качать
Википедию целиком, помечая все адреса как запрещённые. Причина оказалась в том,
что \texttt{RobotFileParser.read()} ходит за \texttt{robots.txt} через
\texttt{urllib} со стандартным User-Agent, Викимедиа отвечает на него
\texttt{403}, а парсер трактует \texttt{403} как \enquote{запрещено всё}.
Исправление --- качать \texttt{robots.txt} той же сессией, что и остальные
страницы, и разбирать его через \texttt{parse()}.

\textbf{Адрес нужно нормализовать раньше, чем кажется.} Ссылки в HTML Википедии
процентно закодированы, и пока адреса не раскрывались, правила с кириллицей по
ним не срабатывали: в базу лезли служебные страницы, которые фильтр должен был
отсечь.

\subsection*{Недостатки}

Робот однопоточный. Пауза вежливости в секунду означает предел около двух
документов в секунду на два источника, и на корпус в 30 тысяч документов уйдут
часы. Чередование источников частично это скрадывает, но настоящее решение ---
несколько потоков с очередью на каждый хост, как в схеме URL frontier с лекции.

\enquote{Сырой} HTML хранится без сжатия, и на него уходит на порядок больше
места, чем на текст. Текст сильно избыточен и сжимается в разы, так что сжатие
здесь напрашивается.

Переобкачка обходит документы по времени последней проверки, то есть одинаково
часто для всех. Разумнее было бы учитывать, как часто документ менялся раньше, и
проверять стабильные страницы реже, а живые --- чаще.

Наконец, робот не защищён от ловушек: календари и страницы с бесконечной
пагинацией порождают бесконечную очередь. Пока от этого спасают только
ограничение глубины и фильтры адресов.

\pagebreak
